% =============================================================================
% Online Appendix to Section 5 (Energy Security): further results.
% Numbers: paper/results/es.json (scripts/paper/fig_es.R, final rerun of 2026-10-12).
% All statistics over rankable cells (config/rankability.yml) unless stated.
% =============================================================================
\section{Energy Security: further results}\label{oa:es}

\paragraph{Rankability.} The rankability rule keeps 4{,}163 of the 6{,}232 scored cells. It removes
the 607 cells of the 24 territories and 1{,}462 cells of sovereign economies whose score rests
mainly on imputation. Among the rankable cells the mean imputed weight share falls to 0.18
upstream, 0.08 midstream and 0.22 downstream. A pair holds between 68 rankable economies (Nuclear
downstream) and 195 (Nuclear upstream). Within each pair every scored economy observes the same
set of categories. The upstream trade baskets give Trade categories to Electric Vehicles, Wind,
Solar and Nuclear upstream, so nine of the ten upstream pairs carry one; Green Hydrogen upstream is
the exception.

\paragraph{Distribution by pair.} In Electric Grid, Solar and Green Hydrogen upstream the
interquartile range is narrower than 0.05. Upstream pair medians run from 0.17 in nuclear and 0.20
in batteries and electric vehicles to 0.47 in oil. Midstream and downstream the pillar rests mainly
on median-centred S-curves; the four midstream pairs scored on Trade alone (gas, coal, oil and
nuclear) have medians between 0.48 and 0.56. Green Hydrogen midstream, scored on Foreign Dependency
alone, gives 90 per cent of its 106 rankable economies the value 0.26; Green Hydrogen downstream,
scored on Consumption alone, is bimodal (interquartile range 0.10 to 0.72). Within stages,
technology-stage effects explain 33, 50 and 22 per cent of the upstream, midstream and downstream
variance, and country effects 51, 65 and 29 per cent of the within-pair remainder. Subtracting
the contributions of Energy Prices and of the world-concentration terms of Trade and Minerals
Trade leaves the pair-effect $R^2$ at 48 per cent.

\paragraph{Drivers.} Investment, defined midstream only for batteries, electric vehicles, solar and
wind, supplies 1.4 per cent of the midstream score against a mean effective weight of 7.0 per cent,
because the investment monitor records neither manufacturing investment nor capacity for most
economies (72 per cent of rankable economies score zero in the median of these pairs). Downstream,
resource potential supplies 10 per cent and Investment 9.8 per cent against a weight of 19 per
cent. The concentration term of Minerals Trade is constant for all economies except Taiwan and
R\'eunion, which lie outside the World Bank list. Energy Prices scores range from 0.19 for gas to
0.81 for the electric grid, raising every economy's upstream score by between 0.025 (batteries)
and 0.088 (nuclear) without changing any within-pair rank.

\paragraph{Correlation among categories.} The median within-pair Spearman correlation between two
categories is 0.42 upstream (one of 78 pair-specific correlations is negative, $-0.01$), 0.33
midstream and 0.11 downstream. Reserves and Production correlate at a median of 0.84 over the
seven pairs that carry both, and at 0.87 among the economies with a positive score on both (a
median of 16 per pair). Foreign Dependency correlates with Production at 0.76 and with Reserves at
0.69. Midstream, Foreign Dependency and Trade correlate at only 0.28 (range $-0.23$ to 0.48):
economies with large shares of refined supply and manufacturing capacity are not always those with
large export shares and small net imports. Downstream, Energy Access and Consumption correlate at
0.12 and resource potential correlates negatively with Consumption ($-0.11$). Within a pair the
pillar's ranking follows Energy Imports (median Spearman correlation with $\ES_{its}$ of 0.80),
Minerals Trade (0.79) and Reserves (0.79) upstream, Trade (0.94) midstream, and Investment (0.76)
and Consumption (0.74) downstream.

\paragraph{Uranium.} An economy with neither resources nor production scores the Energy Prices
contribution, $0.704/8=0.088$, plus its Trade category. For the 105 economies whose uranium trade is
not observed, Trade is imputed and the score is 0.172 (shared rank 67); the other 32 economies
without uranium score between 0.145 (Trinidad and Tobago) and 0.201 (Norway). Their imputed weight
shares lie between one quarter and one half, so all remain rankable. Producers dominate the top,
though not all of it: T\"urkiye, which holds resources but reports no mine output, ranks 18th
(0.34), and the lowest-ranked producer is 26th. Nine of the 15 highest-ranked economies are
coalition members. Among the 26 members without identified resources are France, the United
Kingdom, South Korea, the Netherlands, Norway and Poland.

\paragraph{Leaders.} China's eleven first places are Batteries, Electric Vehicles and Wind upstream,
where its shares of world mineral supply, production and trade dominate; Batteries, Electric Grid,
Electric Vehicles, Solar and Wind midstream; and Batteries, Nuclear and Wind downstream. South
Africa leads green hydrogen upstream, whose mineral weights are dominated by the platinum-group
metals, and Chile the grid and solar upstream, in which copper carries most of the demand weight;
in Electric Grid upstream, whose basket holds copper and aluminium inputs, the Democratic Republic
of the Congo and Peru follow Chile. The midstream pairs scored on Trade alone are led by Poland
(coal), South Korea (gas and oil) and Russia (nuclear, whose basket holds enriched uranium, reactor
components and steam-cycle equipment). Indonesia leads Green Hydrogen midstream, in which 95 of
106 rankable economies share one value. Downstream, Qatar leads oil, Kuwait gas, Indonesia coal,
Saudi Arabia solar, the United States geothermal and green hydrogen, and T\"urkiye electric
vehicles, a pair that rests on Investment alone.

In Oil upstream (Figure~\ref{fig:es_top}) Saudi Arabia leads (0.81); the United States is 12th
(0.71) and China 34th (0.55). The cell-constant inputs give every economy 0.16; of the gap between
the top-ten mean and the pair mean, Reserves accounts for 58 per cent, Energy Imports for 28 and
the country-varying part of Trade for 13. China's Reserves and Trade scores (0.64 and 0.79) are
close to those of the United States (0.69 and 0.87), but its Energy Imports score is 0.25 against
0.68. In Batteries and Solar midstream China leads with 0.77 (United States 0.56) and 0.82 (India
0.67), and coalition members fill the other nine top-ten places of each, among them Indonesia and
Malaysia in batteries and Thailand, Malaysia, Viet Nam and Indonesia in solar. Investment accounts
for 62 and 58 per cent of the top-ten gap in these pairs (mean scores 0.13 and 0.09, top ten 0.61
and 0.58); China's lead over the runner-up comes from Foreign Dependency (0.37 against 0.21 in
batteries, 0.40 against 0.28 in solar), while on Investment the United States scores higher than
China in batteries (0.86 against 0.80).

\paragraph{The coalition's tiers.} The partners, several of which enter as major suppliers, have the
higher upstream mean percentile (0.63 against 0.56 for the core), while the core is ahead midstream
(0.78 against 0.64) and downstream (0.64 against 0.59). The 21 EU members of the coalition are
weaker upstream than the other members (0.50 against 0.64) and stronger midstream (0.79 against
0.67). By mean within-pair percentile the United States (0.95), Canada (0.89), Australia (0.82),
India (0.79) and Poland (0.77) lead, and Qatar (0.44), Ghana (0.45), Gabon (0.46), Namibia (0.46)
and Luxembourg (0.47) trail. Midstream, Japan has the highest mean percentile of any member (0.95),
ahead of Germany, the United States, India and South Korea. Of the 1{,}710 member cells in the 30
scored pairs, 1{,}606 are rankable, 96 scored but not rankable and 8 unscored; the non-rankable
cells concentrate in Namibia, Costa Rica, the Dominican Republic, Guatemala and Ghana (rankable in
13 to 15 of their 28 or 29 scored pairs) and Canada (23 of 29).

\paragraph{The United States, China and Taiwan.} The United States leads two pairs (Geothermal and
Green Hydrogen downstream); its weakest positions are 30th in Nuclear midstream, 23rd in Oil
downstream, 22nd in Electric Grid midstream and 17th in Nuclear downstream. China is in the top ten
of every midstream pair. Its weak positions are those of a fossil-fuel importer (33rd in gas and
34th in oil upstream, 24th in gas downstream) and two downstream pairs in which deployment data
place it low: 63rd in green hydrogen and last of 186 in oil, where Consumption carries five-ninths
of the weight and no rankable economy has lower projected growth of oil-fired capacity. The United
States ranks above China in 8 of the 30 pairs, all upstream or downstream fossil-fuel and resource
pairs and none midstream. Taiwan, scored in 28 cells and rankable in 19, has a midstream mean
percentile of 0.75 (ranks 27th to 61st on Trade) and 0.01 across its seven rankable upstream pairs,
with no place in any top ten.

\paragraph{Ties.} In Green Hydrogen midstream the United States' rank of 12 is shared with 94 other
economies, and in Green Hydrogen downstream 49 economies share third place at 0.72.
